% Accepted after the coordinating agent's independent re-proof; integrated below.
\begin{claim}{C.28}{Proved}{Neither a flat interval nor countable scalar descent is automatic}\label{C28}
There are normalized data satisfying all four manuscript hypotheses and all the corollary's small-projective-generator hypotheses, with a pointwise cocontinuous symmetric tensor, for which no flat object of $\mathcal M$ admits a monomorphism from the unit. Moreover, $\mathcal M$ admits no form over any countable subring of $k$, even up to $k$-linear equivalence of abelian module categories. Thus neither the flat-interval hypothesis of Claim~\ref{C27} nor the countable-form hypothesis of Claim~\ref{C22} follows from the manuscript, even after changing the small linear presentation. All raw resolving adjunctions in this example nevertheless satisfy the desired conclusion by Claims~\ref{C7} and~\ref{C17}; it is not a counterexample to the remaining raw-adjoint gate.
\end{claim}
\begin{proof}
Choose an uncountable set $I$ and the free Boolean ring
\[
 k=\mathbb F_2[e_i\mid i\in I]/(e_i^2-e_i\mid i\in I).
\]
Every element involves only finitely many generators, and every assignment of zeroes and ones to the $e_i$ defines a character to $\mathbb F_2$. Put
\[
 H_i=k\oplus ke_i x_i,\qquad x_i^2=0.
\]
This is a cocommutative bialgebra, with augmentation $\varepsilon_i(x_i)=0$, $1$ grouplike, and $x_i$ primitive. The relation $(1-e_i)x_i=0$ is respected by the coproduct, and the square-zero relation is respected because the characteristic is two. Its underlying $k$-module is finitely generated projective.

Let $\mathcal A$ have objects $u,b_0,b_1,a_i$ for $i\in I$. The endomorphism rings at $u,b_0,b_1$ are $k$, and that at $a_i$ is $H_i$. For each branch $v$ there is one copy $kg_v$ of morphisms $u\to v$; all other off-diagonal Hom modules vanish. Composition with a branch endomorphism is $h g_v=\varepsilon_v(h)g_v$, where $\varepsilon_{b_j}=\id_k$. These rules define a small $k$-linear category with projective Hom modules. Give $g_v$ its grouplike coalgebra structure. Composition preserves these structures, so the pointwise tensor on
\[
 \mathcal M_0=\operatorname{Fun}_k(\mathcal A^{\mathrm{op}},\Mod k)
\]
is well-defined. Explicitly, a presheaf consists of modules $X_u,X_{b_j},X_i$, with $H_i$ acting on $X_i$, and maps $f_v:X_v\to X_u$ satisfying $f_i x_i=0$. On a tensor, $x_i$ acts as $x_i\otimes1+1\otimes x_i$, and the branch map is the tensor of the two branch maps. The usual module symmetry, associativity, and unit maps respect these formulas. Thus the tensor is symmetric, $k$-bilinear, and cocontinuous.

The unit $U$ has value $k$ at every vertex, identity branch maps, and zero $x_i$-actions. Its scalar copower $\kappa_0V$ is the constant presheaf with value $V$. A map into it is determined by its component at $u$, and all branch components are forced to be its composites with $f_v$. Consequently
\[
 \operatorname{Hom}_{\mathcal M_0}(X,\kappa_0V)
 \cong\operatorname{Hom}_k(X_u,V).
\]
This fixes the variance: these are contravariant presheaves on the star. The corresponding constant covariant module is representable at $u$; we are not asserting that the unit presheaf is representable.

Now set $\mathcal M=s\mathcal M_0$, equivalently the presheaf category on $\mathcal A\otimes_k k[\Delta]$, with tensor pointwise in both variables. We use these original indexing objects for the pointwise formula and pass to the additive envelope only for the corollary's additive-generator convention. The category has all small limits and colimits, exact products, and finitely presented projective representable generators. Their Hom modules are finite sums of projective $k$-modules, since the Hom sets of $\Delta$ are finite. Its scalar copower $\kappa(V)$ is constant in both variables, hence is exact and preserves all limits. Its left adjoint and a resolving functor are
\[
 \varphi(X)=H_0N(X_u),\qquad F(X)=N(X_u).
\]
Here $N$ denotes normalized chains. Each $F_n$ is a retract of the single evaluation at $(u,[n])$ by the normalized/degenerate decomposition. On the representable at $(a,[m])$, evaluation at $u$ is $k[\Delta^m]$, because $\mathcal A(u,a)=k$. Its normalized complex has finite free components, vanishes in positive homology, and has $H_0=k$. Claim~\ref{C1} shows that $F$ takes projectives to projective complexes, so it is a projective resolving functor as in Claim~\ref{C6}. Its augmentation on the unit is the identity of $k[0]$, giving the canonical normalization.

Take the simplicial module $L=k[\Delta^1]/k[\{0\}]$, let $J=U\otimes_k L$, and let $\alpha:U\to J$ be induced by the vertex $1$. In every degree the constant-one simplex is a basis element of $L$, so $\alpha\otimes X$ is evaluationwise split monic for every $X$. The pointed simplicial contraction $\min:\Delta^1\times\Delta^1\to\Delta^1$ contracts $L$ to zero. Tensoring this homotopy with $X_u$ and applying the normalized prism homotopy gives
\[
 F(J\otimes X)=N(L\otimes X_u)\simeq0
\]
for all $X$, including degree-zero homology. This verifies (a)--(d). The evaluationwise splittings pass through every retract of a sum of evaluation functors, so the alternative hypothesis of Claim~\ref{C7} holds for every chosen resolving functor. The particular $F$ above also satisfies Claim~\ref{C17}.

We next prove the flat-object obstruction. Let $\mathcal B$ be the full three-vertex subcategory on $u,b_0,b_1$, and let $U_{\mathcal B}$ be its constant presheaf. Write $h_v=\mathcal B(-,v)$ and $h^v=\mathcal B(v,-)$. The covariant module $T$ with $T(u)=k$ and $T(b_j)=0$ has the projective resolution
\[
 0\longrightarrow h^{b_0}\oplus h^{b_1}\longrightarrow h^u
 \longrightarrow T\longrightarrow0.
\]
Tensoring it with $U_{\mathcal B}$ gives the addition map $k^2\to k$, so
\[
 \operatorname{Tor}^{\mathcal B}_1(T,U_{\mathcal B})=k,
 \qquad \operatorname{Tor}^{\mathcal B}_2(T,-)=0.
\]
There cannot be a monomorphism $U_{\mathcal B}\to Q$ with $Q$ flat: the Tor sequence of its cokernel would inject this nonzero $k$ into $\operatorname{Tor}^{\mathcal B}_1(T,Q)=0$.

Restriction to $\mathcal B$ followed by simplicial evaluation in degree zero is exact and preserves filtered colimits. It takes every representable of $\mathcal M$ to a finite sum of representables of $\mathcal B$: the restriction of $h_{a_i}$ is $h_u$, and evaluation of $k[\Delta^m]$ in degree zero has $m+1$ generators. Hence it preserves projectives and flat objects, using the standard description of flat functors as filtered colimits of finitely generated projectives; see \citet[p.~5]{Krause2016C}. It sends the unit to $U_{\mathcal B}$. A monomorphism from the unit into any flat object of $\mathcal M$ would therefore contradict the preceding paragraph. In particular no possible interval is flat. The two finite star matrices are checked in \texttt{code/C\_star\_obstruction.py}; the Tor argument proves the assertion for arbitrary flat objects.

Finally suppose, toward a contradiction, that a countable unital subring $s\subseteq k$ and a small $s$-linear category $\mathcal D_s$ give a $k$-linear abelian equivalence
\[
 \mathcal M\simeq
 \operatorname{Fun}_k\bigl((k\otimes_s\mathcal D_s)^{\mathrm{op}},\Mod k\bigr).
\]
No preservation of the displayed generators or tensor is assumed. All elements of $s$ together involve a countable subset $I_s\subset I$. Let $\chi_0:k\to\mathbb F_2$ send every $e_i$ to zero, and let $\chi_1$ send $e_i$ to zero on $I_s$ and to one outside $I_s$. The characters agree on $s$.

For a $k$-linear abelian category, the full subcategory of objects annihilated by $\ker\chi$ is intrinsic: each scalar in that ideal acts by zero on the identity of each object. A $k$-linear equivalence preserves it. In the proposed scalar form the two fiber subcategories are both module categories over $\mathbb F_2\otimes_s\mathcal D_s$, so they would be equivalent. In our construction they are $s\Mod(\mathcal A_{\chi_j}^{\mathrm{op}})$. At $\chi_0$ every branch algebra is $\mathbb F_2$; at $\chi_1$ every branch outside $I_s$ has algebra $H=\mathbb F_2[x]/(x^2)$.

Every simple base presheaf is supported at one vertex. Indeed its $u$-component defines a subobject supported at $u$; if it is zero, each branch defines a subobject supported at that branch. Each branch algebra has its unique simple module $\mathbb F_2$. At $\chi_0$ the simple at $u$ is projective, and a branch simple has resolution
\[
 0\longrightarrow h_u\longrightarrow h_v\longrightarrow S_v\longrightarrow0.
\]
Thus every simple has finite projective dimension. At an active branch of $\chi_1$, evaluation is exact and preserves projectives: the only nonzero representable evaluation there is $H$. Its simple evaluates to $\mathbb F_2$, which has infinite projective dimension over $H$. The periodic resolution has all terms $H$ and all differentials multiplication by $x$; its image and kernel are both $(x)$. Applying $\operatorname{Hom}_H(-,\mathbb F_2)$ gives zero differentials and a nonzero Ext group in every degree. Consequently the branch simple has infinite projective dimension in the base presheaf category.

Dold--Kan identifies simplicial objects with nonnegative chain complexes. Their simple objects are exactly the spheres $S[n]$ on base simples. Degree evaluation is exact and takes projective complexes to projective objects, so infinite projective dimension of $S$ implies the same for $S[n]$. If $S$ has finite projective dimension $d$, its sphere $S[0]$ and every disk $D^n(S)$ have projective dimension at most $d$; the sequences
\[
 0\longrightarrow S[n-1]\longrightarrow D^n(S)
 \longrightarrow S[n]\longrightarrow0
\]
give finite projective dimension, at most $d+n$, for every sphere. Therefore the $\chi_0$ fiber has no simple object of infinite projective dimension, whereas the $\chi_1$ fiber does. An abelian equivalence preserves simplicity, projectivity, and projective dimension, a contradiction.

This excludes a countable scalar form even up to $k$-linear abelian equivalence. Ordinary $k$-presheaf categories have such a form over $\mathbb F_2$ by linearizing their ordinary indexing category; hence the example also lies outside Claim~\ref{C21}. It remains inside Claims~\ref{C7} and~\ref{C17}, as verified above. The finite certificate checks the star resolution and the dual-number differential; the infinite periodic-resolution and Boolean-character arguments are the proofs of the two infinite assertions, not conclusions drawn from a finite truncation.
\end{proof}

